\documentclass{beamer}
\usetheme{Blackboard}

\usepackage[utf8]{inputenc}
\usepackage[T1]{fontenc}
\usepackage{graphicx}
\usepackage{booktabs}

\title{Chalkboard Notes on Sorting}
\subtitle{A lecture-style Beamer starter}
\author{Prof. Dana Okafor}
\institute{Department of Computer Science}
\date{\today}

\begin{document}

\begin{frame}
  \titlepage
\end{frame}

\begin{frame}{Outline}
  \tableofcontents
\end{frame}

\section{Why Sorting Matters}

\begin{frame}{Motivation}
  \begin{itemize}
    \item Sorting is a building block for search, joins, and deduplication
    \item Comparing algorithms teaches asymptotic analysis directly
    \item We will trace three algorithms on the same input
  \end{itemize}
\end{frame}

\section{Algorithms}

\begin{frame}{Insertion Sort}
  \begin{columns}[T]
    \begin{column}{0.5\textwidth}
      \textbf{Idea}
      \begin{itemize}
        \item Grow a sorted prefix one element at a time
        \item Shift larger elements right to make room
      \end{itemize}
    \end{column}
    \begin{column}{0.5\textwidth}
      \textbf{Cost}
      \begin{itemize}
        \item Best case: $O(n)$
        \item Worst case: $O(n^2)$
      \end{itemize}
    \end{column}
  \end{columns}
\end{frame}

\begin{frame}{Merge Sort}
  \begin{block}{Divide and conquer}
    Split the array in half, sort each half recursively, then merge the two
    sorted halves in linear time.
  \end{block}
  \[
    T(n) = 2T(n/2) + O(n) \implies T(n) = O(n \log n)
  \]
\end{frame}

\begin{frame}{Comparing the Three}
  \begin{table}
    \centering
    \begin{tabular}{lccc}
      \toprule
      Algorithm & Best & Average & Worst \\
      \midrule
      Insertion sort & $O(n)$ & $O(n^2)$ & $O(n^2)$ \\
      Merge sort & $O(n\log n)$ & $O(n\log n)$ & $O(n\log n)$ \\
      Quicksort & $O(n\log n)$ & $O(n\log n)$ & $O(n^2)$ \\
      \bottomrule
    \end{tabular}
    \caption{Time complexity by case.}
  \end{table}
\end{frame}

\section{Worked Example}

\begin{frame}{Tracing Merge Sort on [5, 2, 8, 1]}
  \begin{itemize}
    \item Split: [5, 2] and [8, 1]
    \item Split again: [5], [2], [8], [1]
    \item Merge: [2, 5] and [1, 8]
    \item Merge: [1, 2, 5, 8]
  \end{itemize}
\end{frame}

\section{Wrap-up}

\begin{frame}{Summary}
  \begin{itemize}
    \item Insertion sort is simple and fast on nearly-sorted input
    \item Merge sort guarantees $O(n \log n)$ but needs extra memory
    \item Quicksort is usually fastest in practice despite a worse worst case
  \end{itemize}
\end{frame}

\begin{frame}{Questions}
  \centering
  \Large Questions?\\[1em]
  \normalsize dana.okafor@example.edu
\end{frame}

\end{document}
